% コンパイル: tectonic docs/theory_ja.tex
%   （XeLaTeX 系。lualatex を使う場合は下の documentclass を
%    \documentclass[11pt,a4paper]{ltjsarticle} に置き換え、
%    xeCJK のブロックを削除する）
\documentclass[11pt,a4paper]{article}
\usepackage{xeCJK}
\setCJKmainfont{Noto Serif CJK JP}
\setCJKsansfont{Noto Sans CJK JP}
\setCJKmonofont{Noto Sans Mono CJK JP}
\usepackage{amsmath,amssymb,amsthm}
\usepackage{booktabs}
\usepackage{geometry}
\geometry{top=18mm,bottom=20mm,left=17mm,right=17mm}
\usepackage{hyperref}
\hypersetup{colorlinks=true,linkcolor=blue,citecolor=blue,urlcolor=blue}

\newcommand{\R}{\mathbb{R}}
\newcommand{\SOthree}{\mathrm{SO}(3)}
\newcommand{\dockq}{\mathrm{DockQ}}
\newtheorem{prop}{命題}
\newtheorem{defn}{定義}

\title{学習可能なFFTドッキング\\
{\large スコア関数の線形性を使って「どこで損をしているか」を測る}}
\author{DDOCK プロジェクト実験ノート}
\date{2026年8月}

\begin{document}
\maketitle

\begin{abstract}
2つのタンパク質がどう結合するかを予測する剛体ドッキングにおいて、
古典的手法ZDOCKのスコア関数の$144$個のパラメータだけを学習で決め直した。
本稿は、その理論的背景と、実験の経緯を数式とともに整理する。

中心となる観察は、\textbf{スコアが学習パラメータについてアフィンである}という
一点である。これにより、(i)~学習が凸になるので局所解の心配が消え、
(ii)~複合体ごとの分離可能性を線形計画で、共有表の達成値を凸代理問題で診断でき、
(iii)~「損失が悪いのか・容量が足りないのか・最適化が下手なのか」を分離して
測定できる。この性質を用いて誤差を4つの成分へ分解し、
公式ベンチマークで acceptable $66.4\% \to 84.0\%$、high $30.8\% \to 73.6\%$ を得た。
ただし後者の伸びの大半は学習ではなく後処理（局所精密化）によるものである。
同時に、データ量の増加・パラメータ数の増加・hard negative採掘が
\emph{効かない}ことを再現可能な形で示した。
\end{abstract}

\tableofcontents

\section*{用語}
\addcontentsline{toc}{section}{用語}

本文で繰り返し使う語を先にまとめる。

\begin{description}\setlength{\itemsep}{0pt}
  \item[PSC] pairwise shape complementarity。形状相補性。
        表面どうしの重なりを加点し、内部どうしの重なりを減点する項。
  \item[IFACE] 原子タイプ対の統計ポテンシャル。本研究で学習する$144$個の表そのもの。
  \item[ELEC] 静電相互作用項。
  \item[pose] リガンドの1つの配置$g \in \mathrm{SE}(3)$。
  \item[holo / apo] holoは結合型（複合体から取り出した構造）、
        apoは非結合型（単独で決定した構造）。holo同士の再ドッキングが最も易しい。
  \item[固定プール] 一度だけ探索を回して書き出した候補poseの集合を固定し、
        以後は表を変えても\emph{並べ替えだけ}を行う評価法。安価だが、
        表を変えると本来は探索が返す候補も変わるので、end-to-endとは別物である。
  \item[end-to-end] 学習した表で\emph{探索からやり直す}評価。実運用に対応する。
  \item[候補回収 (retrieval)] 正解に近いposeが候補集合に\emph{含まれているか}。
        それを選べるかは別問題（順位づけ）である。
  \item[seed] 乱数の初期値。分割・初期化・採掘順などを決める。
  \item[AUC] ROC曲線下面積。$0.5$が当てずっぽう、$1.0$が完全な判別。
  \item[NMS] non-maximum suppression。近すぎるposeを間引いて多様性を保つ操作。
  \item[InfoNCE] 正例の確率質量を負例に対して持ち上げる対照学習の損失。
  \item[softplus] $\mathrm{softplus}(x) = \log(1+e^x)$。$\max(0,x)$の滑らかな近似。
  \item[LSE] log-sum-exp、$\mathrm{LSE}(x) = \log\sum_i e^{x_i}$。$\max$の滑らかな近似。
  \item[big-$M$] 整数計画で「この制約は無効にしてよい」を表すための十分大きい定数。
  \item[L-BFGS] 準ニュートン法の一種。凸で滑らかな問題を少ない反復で解く。
  \item[pp] percentage point。$30\% \to 40\%$は$+10$\,pp（$+33\%$ではない）。
\end{description}

\section{問題設定}

\subsection{剛体ドッキング}

2つのタンパク質$A, B$の立体構造が別々に既知であるとき、
両者が結合した複合体の構造を予測する問題を考える。
$A$を\textbf{受容体}（固定）、$B$を\textbf{リガンド}（可動）と呼ぶ。
タンパク質を剛体とみなすと、$B$の配置は
\begin{equation}
  g = (R, t) \in \mathrm{SE}(3) = \SOthree \ltimes \R^3
\end{equation}
という$6$自由度で表される。$R$は回転、$t$は並進である
（$\ltimes$は半直積。並進と回転は可換でないので直積ではない）。
リガンドの原子座標を$\{y_1,\dots,y_{N_L}\}$とすると、配置$g$における座標は
\begin{equation}
  y_i(g) = R\,(y_i - \bar y) + \bar y + t
\end{equation}
となる（$\bar y$は重心）。この$g$を\textbf{pose}と呼ぶ。

予測の目標は、正解の配置$g^\star$（実験構造から定まる）に近い$g$を見つけることである。

\subsection{スコア関数}

配置の良し悪しを実数値で測る関数$S(g)$を\textbf{スコア関数}と呼び、
$\arg\max_g S(g)$を答えとする。ZDOCK 3.0.2 のスコアは3項からなる：
\begin{equation}
  S(g) \;=\; \underbrace{\alpha\,S_{\mathrm{PSC}}(g)}_{\text{形状相補性}}
        \;-\; \underbrace{\sum_{i=1}^{12}\sum_{j=1}^{12} e_{ji}\,T_{ij}(g)}_{\text{原子対の統計ポテンシャル}}
        \;+\; \underbrace{\beta\,S_{\mathrm{ELEC}}(g)}_{\text{静電相互作用}} .
  \label{eq:score}
\end{equation}

ここで中央の項が本研究の対象である。
全原子を$12$種類の\textbf{原子タイプ}に分類し
（実際には $(\text{残基名},\text{原子名})$ の$167$通りの組を$12$個へ写像する）、
接触特徴$T_{ij}(g)$を数える。実装は原子対を直接数えるのではなく、格子を経由する：
\begin{equation}
  T_{ij}(g) \;=\; \sum_{v \in \Lambda} L_i(v; g)\, H_j(v),
  \label{eq:tij}
\end{equation}
ここで$L_i(v;g)$は「配置$g$でタイプ$i$のリガンド原子が格子セル$v$を占めるか」の指示関数、
$H_j(v)$は「格子点$v$から$r_c = 6$\,\AA{}以内にあるタイプ$j$の受容体原子の個数」である。
すなわち\textbf{第1添字がリガンド、第2添字が受容体}であり、
$T_{ij}$は「タイプ$i$のリガンド原子とタイプ$j$の受容体原子の接触量」を測る。

$e_{ij}$がその重みで、$12\times12 = 144$個ある。
式\eqref{eq:score}に転置と負号が付いているのは実装の約束による。
公開表$e$は「(受容体タイプ, リガンドタイプ)」の順で格納され、しかも
「有利な接触ほど負」という符号規約なので、スコアが実際に$T$と縮約する行列は
\begin{equation}
  W(e) = \mathrm{IFACE\_SIGN}\cdot e^{\top} = -\,e^{\top}
  \qquad (\texttt{src/zdock/score.py:491})
\end{equation}
である。これを並べたベクトルを
\begin{equation}
  e = (e_{11},\dots,e_{12,12}) \in \R^{144}
\end{equation}
と書く。\textbf{本研究で学習するのはこの$e$だけ}であり、
$\alpha, \beta$ および$S_{\mathrm{PSC}}$内部の係数は固定する。

\subsection{FFTによる全探索}

$6$次元空間$\mathrm{SE}(3)$を素朴に離散化すると組合せ爆発する。
ZDOCKは以下の構造を使って回避する。

まず$\SOthree$を有限集合
$\mathcal{Q} = \{R_1,\dots,R_M\}$（$M = 1944$）に離散化する。
各$R_m$を固定すると、並進についてのスコアは\textbf{相関}の形になる。
受容体とリガンドをそれぞれ格子上の関数
$F^{(k)}_R : \Lambda \to \R$、$F^{(k)}_L : \Lambda \to \R$
（$\Lambda$は$1.2$\,\AA{}間隔の3次元格子、$k$はチャネル）で表すと、
\begin{equation}
  T_{ij}(R_m, t) \;=\; \sum_{u \in \Lambda} F^{(i)}_R(u)\, F^{(j)}_L(u - t)
  \;=\; \bigl(F^{(i)}_R \star F^{(j)}_L\bigr)(t)
  \label{eq:corr}
\end{equation}
と書ける。相関は畳み込み定理により
\begin{equation}
  F^{(i)}_R \star F^{(j)}_L
  = \mathcal{F}^{-1}\Bigl[\mathcal{F}[F^{(i)}_R] \cdot \overline{\mathcal{F}[F^{(j)}_L]}\Bigr]
  \label{eq:fft}
\end{equation}
と計算でき、格子点数$n$に対して$O(n\log n)$で\textbf{全並進位置}が一度に得られる。
共役が付くのは\textbf{リガンド側}である（$\overline{\;\cdot\;}$は複素共役）。
式\eqref{eq:corr}の定義で受容体側に共役を付けると$t$の符号が反転してしまう。
実装も同じ約束である（\texttt{src/zdock/search.py}）。

\begin{prop}[FFT計算可能性の十分条件]
固定した回転$R_m$について全並進を一度のFFTで評価できるのは、スコアが
\emph{有限個のチャネル相関の和}
\begin{equation}
  S(R_m, t) = \sum_{k} \bigl(F^{(k)}_R \star F^{(k)}_{L,R_m}\bigr)(t)
\end{equation}
に分解できるときである。原子対についての和
\begin{equation}
  S(g) = \sum_{a \in A}\sum_{b \in B} w\bigl(\mathrm{type}(a), \mathrm{type}(b), \|x_a - y_b(g)\|\bigr)
\end{equation}
という形（\textbf{ペア加法的}）は、この条件を満たす代表的な構成である。
\end{prop}

\textbf{必要条件ではない}点に注意する。相関の非線形な組合せなど、
ペア加法でなくてもFFT後に全並進を評価できる関数は存在する。
ここで主張しているのは「ペア加法なら必ずFFTに乗る」という一方向だけである。

ただし本実装が採用しているのはペア加法の構成であり、その制約は本研究を通じて
繰り返し現れる。原子タイプを細分しても距離を多段にしても加法性は保たれるが、
「この接触は、あの接触も同時にあるときだけ有利」という協同効果は表現できない。

\textbf{PSC項だけは例外的な形をとる。} 形状相補性は受容体・リガンドを複素格子
$Z_R, Z_L$（実部が表面層、虚部がコア）で表し、その積の実部を取るので、
実数相関ではなく
$\mathcal{F}^{-1}\bigl[\mathcal{F}[Z_R]\cdot\overline{\mathcal{F}[\overline{Z_L}]}\bigr]$
の実部として計算される。$O(n\log n)$である点は変わらない。

各回転について並進の上位を取り、全体で上位$K$件（既定$K = 1500$）を残す。

\section{評価}

\subsection{DockQとCAPRI分類}

予測poseと正解の一致度は$\dockq \in [0,1]$で測る。
これは接触の再現率$f_{\mathrm{nat}}$、界面RMSD、リガンドRMSDを合成した量で、
国際コンペCAPRIの基準により4段階に離散化される：
\begin{equation}
  \mathrm{CAPRI}(\dockq) =
  \begin{cases}
    \text{Incorrect}  & \dockq < 0.23,\\
    \text{Acceptable} & 0.23 \le \dockq < 0.49,\\
    \text{Medium}     & 0.49 \le \dockq < 0.80,\\
    \text{High}       & \dockq \ge 0.80 .
  \end{cases}
\end{equation}

\subsection{報告する指標}

複合体の集合$\mathcal{C}$（$|\mathcal{C}| = 250$）について、
各複合体で上位$K$件を提出したとき
\begin{align}
  \mathrm{Max}(\mathrm{Top}\,k)_{\text{acc}}
  &= \frac{1}{|\mathcal{C}|}\sum_{c \in \mathcal{C}}
     \mathbf{1}\Bigl[\max_{p \in \mathrm{Top}_k(c)} \dockq(p) \ge 0.23\Bigr],\\
  \mathrm{Oracle}@K &= \mathrm{Max}(\mathrm{Top}\,K) .
\end{align}
$\mathrm{Max}(\mathrm{Top}\,1)$が実運用の指標であり、
$\mathrm{Oracle}@K$は「候補の中から常に正しく選べたら」という上限である。
両者の差が\textbf{順位づけの損失}を表す。

\section{学習問題}

\subsection{アフィン性}

式\eqref{eq:score}を$e$について整理する。
$\alpha, \beta$ および$S_{\mathrm{PSC}}$の係数が固定であるから、
pose $p$のスコアは
\begin{equation}
  \boxed{\;s_p(e) \;=\; c_p \;+\; \langle A_p,\, e\rangle\;}
  \label{eq:affine}
\end{equation}
と書ける。ここで
\begin{align}
  c_p &= \alpha\, S_{\mathrm{PSC}}(p) + \beta\, S_{\mathrm{ELEC}}(p)
         \quad(\text{$e$に依存しない定数}),\\
  A_p &= -\,\mathrm{vec}\bigl(T(p)^{\!\top}\bigr) \in \R^{144}.
\end{align}
負号と転置は$W(e) = -e^\top$から来る。実際
$\sum_{ij}W_{ij}T_{ij} = -\sum_{ij}e_{ji}T_{ij}
= -\langle \mathrm{vec}(e), \mathrm{vec}(T^{\!\top})\rangle$である。
実装で再構成誤差が$0$であることを確認済みである
（\texttt{scripts/capacity\_ceiling.py}の起動時チェック）。

式\eqref{eq:affine}が本研究の出発点である。以下のすべてがここから従う。

\subsection{目的：Top-1の的中}

複合体$c$について、候補集合$\mathcal{P}_c$（正例、$\dockq \ge 0.23$）と
$\mathcal{N}_c$（負例）が与えられる。
「複合体$c$を解けた」とは
\begin{equation}
  \exists\, p \in \mathcal{P}_c \;\text{s.t.}\; \forall n \in \mathcal{N}_c:\;
  s_p(e) > s_n(e)
  \label{eq:solved}
\end{equation}
であり、最大化したいのは
\begin{equation}
  H(e) \;=\; \sum_{c} \mathbf{1}\bigl[\text{式\eqref{eq:solved}が成立}\bigr].
  \label{eq:objective}
\end{equation}

式\eqref{eq:solved}を$e$について書き直すと、
\begin{equation}
  \langle A_p - A_n,\; e\rangle \;>\; c_n - c_p
  \label{eq:linear-ineq}
\end{equation}
という\textbf{線形不等式}になる。$\exists$（論理和）があるため
$H$自体は非凸だが、この構造が後の議論をすべて可能にする。

\subsection{当初の損失関数}

本研究は当初、次の3項を最小化していた（$s$は複合体ごとに中心化し標準偏差で割った値）：
\begin{equation}
  \mathcal{L}(e) \;=\;
  \underbrace{\mathcal{L}_{\mathrm{basin}}}_{\text{(a)}}
  + 0.5\,\underbrace{\mathcal{L}_{\mathrm{margin}}}_{\text{(b)}}
  + 0.1\,\underbrace{\|e - e_0\|^2}_{\text{(c)}} .
  \label{eq:oldloss}
\end{equation}

\paragraph{(a) basin損失（multi-positive InfoNCE）}
離散化と対称性により正例は複数存在するので、
その集合全体の確率質量を持ち上げる：
\begin{equation}
  \mathcal{L}_{\mathrm{basin}}
  = -\log \frac{\sum_{p \in \mathcal{P}} e^{s_p/\tau}}{\sum_{q \in \mathcal{P}\cup\mathcal{N}} e^{s_q/\tau}},
  \qquad \tau = 0.5 .
\end{equation}
分配関数を
\begin{equation}
  Z_{\mathcal{P}} = \sum_{p \in \mathcal{P}} e^{s_p/\tau}, \qquad
  Z_{\mathcal{N}} = \sum_{n \in \mathcal{N}} e^{s_n/\tau}
\end{equation}
と置くと、分母が$Z_{\mathcal{P}} + Z_{\mathcal{N}}$なので
\begin{equation}
  \mathcal{L}_{\mathrm{basin}}
  = \log\Bigl(1 + \frac{Z_{\mathcal{N}}}{Z_{\mathcal{P}}}\Bigr)
  = \mathrm{softplus}\bigl[\log Z_{\mathcal{N}} - \log Z_{\mathcal{P}}\bigr]
\end{equation}
と変形できる（$\mathrm{softplus}(x) = \log(1+e^x)$）。
$\log Z$ は log-sum-exp（LSE）であり、これは$\max$の\textbf{滑らかな近似}である：
$\max_i x_i \le \mathrm{LSE}(x) \le \max_i x_i + \log(\text{要素数})$。
したがって$\mathcal{L}_{\mathrm{basin}}$は実質的に
「最良の正例 vs.\ 最良の負例」を滑らかに比較しており、
$\tau$が小さいほど最良の要素が支配する。

\paragraph{(b) hard negative損失}
\begin{equation}
  \mathcal{L}_{\mathrm{margin}}
  = \frac{1}{|\mathcal{N}|}\sum_{n \in \mathcal{N}}
    \max\Bigl(0,\; m + s_n - \min_{p \in \mathcal{P}} s_p\Bigr), \qquad m = 1 .
  \label{eq:minanchor}
\end{equation}

\paragraph{(c) 事前分布}
公開表$e_0$からの逸脱を罰する二次項。

\section{診断I：損失の欠陥}

\subsection{アンカーの位置}

式\eqref{eq:minanchor}は$\min_{p} s_p$、すなわち\textbf{最悪の正例}を基準にする。
しかしTop-1の的中に必要なのは$\max_p s_p > \max_n s_n$だけである。
実測すると（複合体ごとの標準偏差を単位として）
\begin{equation}
  \operatorname{median}_c\Bigl[\max_{p} s_p - \min_{p} s_p\Bigr] = 7.10\ \mathrm{SD}
\end{equation}
であり、マージン$m = 1$と合わせて\textbf{約8\,SD分}だけ過剰に厳しい条件を課していた。

その帰結は「厳しすぎる」に留まらない。この項が有効（$\max(0,\cdot)$の中が正）
となる負例の割合は、$\min$アンカーで中央値$74.95\%$、
$\max$アンカーなら中央値$0\%$である（比にして約$740$倍）。
$|\mathcal{N}| \approx 1494$で割っているので、
\textbf{勝敗を決める上位数件（$\max_p s_p$を上回るのは中央値3件）の勾配が
約1000倍に希釈}されていた。名前に反して、これは
hard negative項ではなく\emph{負例全体を押し下げる広域正則化}であった。

\subsection{事前分布による到達不能性}

学習後の表は$\|e - e_0\|_2 = 1.37$の位置にあった。一方、
後述する凸最適化がCVで選ぶ解は$6.61$の位置にある。
罰則$0.1\|e-e_0\|^2$の値はそれぞれ
\begin{equation}
  0.1 \times 1.37^2 = 0.188, \qquad 0.1 \times 6.61^2 = 4.37
\end{equation}
であり、後者はデータ項全体（実測$\approx 2.95$）より大きい。
\textbf{良い解のある領域へ行くことを、旧目的関数は強く不利にしていた}
（厳密に禁じていたわけではない。罰則は有限なので、
データ項が十分大きければ到達しうる。実際には到達しなかった）。

さらに、checkpointを選ぶvalidation損失にもこの罰則が含まれていた。
罰則は検証データに依存しないので、選択基準は
「予測が良い表」ではなく「公開表に近い表」を測っていたことになる。

\section{診断II：凸最適化と容量の上限}

\subsection{凸な代理問題}

式\eqref{eq:solved}の$\exists$を、複合体ごとに1つの正例$p(c)$を
\textbf{固定}することで除去すると、目的関数は凸になる：
\begin{equation}
  \min_{e}\;
  \frac{1}{|\mathcal{C}|}\sum_{c}
  \max\Bigl(0,\; \max_{n \in \mathcal{N}_c}\bigl[\epsilon + s_n(e) - s_{p(c)}(e)\bigr]\Bigr)
  + \frac{\lambda}{2}\|e - e_0\|^2 .
  \label{eq:convex}
\end{equation}
$\max$の中は$e$のアフィン関数なので、その$\max$は凸、$\max(0,\cdot)$も凸、
和と二次項を足しても凸である。したがって
\textbf{局所最小はすべて大域最小であり、悪い局所解に捕まることがない}。
学習率・初期値・停止規則・checkpoint選択という不確定要素がすべて消える。

\textbf{ただし「大域最適が得られた」と断言はできない}。
式\eqref{eq:convex}は$\max$を含む非滑らかな関数であり、
実装はL-BFGSによる数値解である。凸性が保証するのは
「十分に収束した点は大域最適」までで、数値出力そのものが最適性の証明ではない。
消えたのは学習率とepoch選択であって、solverの許容誤差と収束確認は残る。

$\lambda$は交差検証で選ぶ。実測では$\lambda$を4桁動かしても
CV性能は$49.14\%$--$49.76\%$とほぼ平坦であり、
\textbf{正則化の強さはレバーではなかった}。効いたのは目的関数の形である。

\subsection{個別の分離可能性と、共有表の達成値}

式\eqref{eq:linear-ineq}が線形不等式であることから、
「複合体$c$を解ける$e$は存在するか」は\textbf{線形計画の実行可能性問題}になる。
validation 250件のうち候補プールにacceptable poseがあるのは$152$件で、
以下はその$152$件を母数とする。

複合体ごとに\emph{独立に}解くと、$\|e-e_0\|_\infty \le 1$の範囲で
$152/152 = 100\%$が解ける（学習済み表と同じ移動幅$0.316$に制限しても$89.5\%$）。

\begin{center}
\begin{tabular}{lr}
\toprule
 & $152$件中の的中 \\
\midrule
学習済み表（旧SGD、fitで学習しvalで採点） & $65 = 42.8\%$ \\
共有表の凸ヒンジ最適解（fitで解きvalで採点） & $71 = 46.7\%$ \\
同（valで解きvalで採点、in-sample） & $86 = 56.6\%$ \\
複合体ごとに専用の表（個別分離可能性） & $152 = 100\%$ \\
\bottomrule
\end{tabular}
\end{center}

\textbf{$100\%$を「$144$個の容量の上限」と読んではいけない。}
これは複合体ごとに\emph{別々の}$e_c$を許した値であり、
共有表がそこまで届く保証はない（複合体Aが上げろと要求する成分を
複合体Bが下げろと要求しうる）。この量が決定的なのは\emph{低いとき}だけである
——低ければ容量不足が確定するが、高くても何も確定しない。
実際に確定するのは「$144$個で個別には足りている」ことだけである。

同様に$46.7\%$も上限ではない。これは特定のアンカーと特定の$\epsilon$を
固定した\textbf{ヒンジ代理}の最適値であり、$0/1$の的中数そのものの最大値ではない。
別のアンカーならこれを超えうる。$46.7 - 42.8 = 3.9$\,pp は
「損失の作り直しで取れる分」の\emph{下界}であって上界ではない。

一方、この表から確実に読めるのは\textbf{汎化が支配的な制約}だということである。
同一の解法が、当てはめた集合そのものでは$56.6\%$、別集合では$46.7\%$
——約$10$\,pp が汎化で失われている。

共有表の真の最大値$H^\star = \max_e H(e)$を求めるには、
式\eqref{eq:objective}の$\exists$を二値変数で表した混合整数計画が要る：
\begin{align}
  \max_{e,\,y,\,z} \quad & \sum_c y_c \\
  \text{s.t.}\quad
  & \sum_{p \in \mathcal{P}_c} z_{cp} = y_c, \qquad \forall c \\
  & \langle A_p - A_n,\, e\rangle \ \ge\ c_n - c_p + \epsilon - M(1 - z_{cp}),
    \quad \forall c, p, n \\
  & \|e - e_0\|_\infty \le R, \qquad y, z \in \{0,1\} .
\end{align}
$M$は十分大きい定数（big-M）で、$z_{cp} = 0$のとき制約を無効化する。
負例の制約は膨大なので、切除平面法（違反したものだけ追加）で扱う。

実測では$n = 12$複合体で厳密解が得られ（$12/12$、ヒンジ解は$11$）、
$n \ge 25$では上界が自明化して破綻した。
また、この最適解を\emph{学習器}として使うと
公開表より汎化が悪い（$28.9\%$ vs.\ $38.8\%$）。
マージンも正則化もなく的中数だけを最大化するため、
「スコア差$0.001$で勝った」複合体を量産し、新しいデータで反転するからである。

\section{診断III：回転格子が品質を制約する}

\subsection{離散化誤差}

正解の回転$R^\star$は格子$\mathcal{Q}$上にない。最近接格子点までの角度を
\begin{equation}
  \theta^\star = \min_{m} d\bigl(R_m, R^\star\bigr)
\end{equation}
とすると（$d$は回転間の測地距離）、実測で中央値$9.5^\circ$、最大$14.4^\circ$である
（別の集合での測定では中央値$9.73^\circ$、p95 $13.65^\circ$。同じ量の独立な2回の測定である）。
\emph{格子点そのものの}角度誤差はスコア関数では消せない。
ただし後述するとおり、これは「最も近い格子点」の話であって、
別の格子点がより良いposeを与えることは十分ありうる。

\subsection{角度と品質の関係}

正解poseをランダムな軸まわりに角度$\theta$だけ回して公式採点した結果：

\begin{center}
\begin{tabular}{lccccccc}
\toprule
$\theta$ & $0^\circ$ & $2^\circ$ & $4^\circ$ & $6^\circ$ & $8^\circ$ & $10^\circ$ & $12^\circ$\\
\midrule
Acceptable & 100 & 100 & 100 & 100 & 100 & 100 & 99.6\\
Medium     & 100 & 100 & 100 & 99.6 & 99.6 & 98.4 & 96.0\\
\textbf{High} & \textbf{100} & \textbf{100} & \textbf{98.8} & \textbf{81.2} & \textbf{48.4} & \textbf{18.0} & \textbf{4.8}\\
$\dockq$中央値 & 1.000 & 0.970 & 0.930 & 0.870 & 0.800 & 0.740 & 0.680\\
\bottomrule
\end{tabular}
\end{center}

\textbf{集合全体でHigh率を$90\%$程度に保つには$\theta \lesssim 5^\circ$が要る}
（個々の複合体の必要条件ではない。$\theta = 6^\circ$でも$81.2\%$はHighである）。
一方Acceptableは$12^\circ$でもほぼ無傷である。
すなわち回転格子は\emph{品質}のみを制約する。

\subsection{「格子の上限」は上限ではなかった（訂正）}

「正解に最も近い格子回転 $+$ 正解に最も合う並進」で作ったposeを$250$件書き出し、
公式採点させると Acceptable $100\%$、Medium $99.6\%$、\textbf{High $39.6\%$} となる。
学習後の実測が$38.0\%$だったので、当初これを
「格子の上限であり、到達可能分の$96\%$を既に取っている」と解釈した。

\textbf{この解釈は誤りである。} 後の測定で、探索が返した上位$50$件の
生の$\mathrm{Oracle}@50$のHighが$\mathbf{44.80\%}$と出た。
これは同じ格子の上を動いているのに$39.6\%$を超えている。
$39.6\%$が上限でないことの直接の反例である。

原因は定義にある。$39.6\%$は\emph{最近接}格子点
$\arg\min_m d(R_m, R^\star)$
で作った\textbf{参照値}であって、$\max_m \dockq$ではない。
$\dockq$は測地距離の単調関数ではなく界面の接触で決まるので、
角度がやや遠い別の格子点の方が良い$\dockq$を出すことがある。
正しくは
\begin{equation}
  \underbrace{39.6\%}_{\text{最近接格子点の参照値}}
  \;\le\; \underbrace{\max_{m}(\cdot)}_{\text{格子上の真の上限}},
\end{equation}
であり、真の上限については$\ge 44.80\%$という実測下界しか分かっていない。

したがって「スコア関数をどう改良してもHighはこれ以上上がらない」とは言えない。
言えるのは、精密化（格子の外へ出る操作）がHigh Oracleを
$44.80 \to 79.60\%$へ動かしたという事実だけである。
これは離散化が重い制約であることを強く示すが、
\emph{原因が格子だけ}であることを証明してはいない。

$\SOthree$の被覆半径は回転数$M$に対しておよそ$M^{-1/3}$で縮むので、
角度を半分にするには$M$を$8$倍にする必要がある。
$\theta$の中央値を$9.7^\circ \to 4.9^\circ$にするには$M: 1944 \to 15552$である。

\subsection{局所精密化}

格子を細かくする代わりに、探索で得たposeを$\mathrm{SE}(3)$上で連続的に動かす。
$6$パラメータ$x = (\omega_1,\omega_2,\omega_3, t_1,t_2,t_3)$による剛体変換
\begin{equation}
  g_x(y) = \exp\bigl([\omega]_\times\bigr)(y - \bar y) + \bar y + t
\end{equation}
（$[\omega]_\times$は歪対称行列、Rodriguesの公式）について
$S$を最大化する。

勾配法を用いない理由がある。$S_{\mathrm{PSC}}$は格子への最近接割り当てと
硬い距離しきい値を含むため、座標について\textbf{区分定数}であり、
勾配がほとんど至るところ$0$である。よってパターン探索
（各軸の$\pm$方向$12$候補を一括評価し、改善がなければ歩幅を半減）を用いる。

事前に、正解の近傍でスコアの極大がどこにあるかを測った：
\begin{center}
\begin{tabular}{lcc}
\toprule
 & 公開表 & 学習後の表 \\
\midrule
極大の回転ずれ（中央値） & $1.4^\circ$ & $\mathbf{1.0^\circ}$ \\
極大の並進ずれ & $1.02$\,\AA & $\mathbf{0.84}$\,\AA \\
そこでの$\dockq$ & $0.849$ & $\mathbf{0.890}$ \\
\bottomrule
\end{tabular}
\end{center}
極大は正解と一致しないが、\textbf{Highが要求する$5^\circ$の内側}にある。
したがって登れば品質は上がる。また学習が
\emph{スコアの地形そのもの}を改善していたことが分かる。

結果（上位$K = 5$を精密化）：High $38.0\% \to 68.4\%$
（$78$勝$2$敗、$p = 5\times10^{-21}$）。
参照値$39.6\%$を大きく超えたのは、精密化が格子の外へ出る操作だからである。

\section{誤差の分解}

上位$50$件を精密化した後の測定：

\begin{center}
\begin{tabular}{lccc}
\toprule
 & Acceptable & Medium & High\\
\midrule
Top-1（実運用） & $84.00$ & $82.40$ & $73.60$\\
$\mathrm{Oracle}@50$（精密化前） & $94.80$ & $90.00$ & $44.80$\\
$\mathrm{Oracle}@50$（精密化後） & $\mathbf{94.80}$ & $91.60$ & $\mathbf{79.60}$\\
\midrule
（参考）HDOCK & $94.4$ & $94.0$ & $93.6$\\
\bottomrule
\end{tabular}
\end{center}

読み方は次のとおり。
\begin{itemize}
  \item \textbf{候補回収}：$\mathrm{Oracle}@50$のAcceptableが$94.80\%$。
        正解は候補に入っている。探索の改良は不要。
  \item \textbf{品質}：精密化前のHigh Oracle $44.80\%$は候補が無いのではなく
        \emph{候補の品質}の問題であり、精密化が$79.60\%$まで回復させた。
        離散化がその主因と考えられるが、上で述べたとおり証明はしていない。
  \item \textbf{順位づけ}：Acceptableで$94.80 - 84.00 = 10.8$\,pp を選び間違えている。
  \item \textbf{残る品質不足}：High Oracle $79.60\%$は
        HDOCKの$93.6\%$に届かない。完璧に選んでも足りない。
\end{itemize}

\section{効かなかったもの（negative results）}

\subsection{データ量}

訓練複合体数$N$を$220 \to 500 \to 1000$と増やしても改善しない。
公式指標で$74.80 \to 72.67 \to 72.93\%$（各3 seed）であり、
$6$組のseed対すべてで$N=220$が勝ち、負けた複合体は一つもない。

凸最適化なら学習率もseedも介在しないので、学習曲線をきれいに引ける。
固定validation $152$件でのTop-1（\texttt{scripts/capacity\_series.py}）：
\begin{center}
\begin{tabular}{lrrrrr}
\toprule
学習複合体数 $n$ & 50 & 100 & 200 & 400 & 637 \\
\midrule
sym $(12)$   & $48.9$ & $48.7$ & $47.8$ & $48.7$ & $50.0$ \\
full $(144)$ & $35.7$ & $41.7$ & $42.8$ & $45.6$ & $46.1$ \\
\bottomrule
\end{tabular}
\end{center}
小さいモデルは$n = 50$で既に飽和しており、
$144$個を決めるにも$n \approx 400$でほぼ止まる。

\textbf{ただしこの曲線の射程は限られる}（この但し書きは重要である）。
曲線は候補プールに正例を持つ$637$件だけで引いており、
部分集合の抽出は3反復しかしていない。
また上の$N = 220/500/1000$の比較は旧SGDのもので、
学習ステップ数との交絡が完全には除けていない。
「データ量は一般に効かない」ではなく
「\emph{この}特徴量・\emph{この}規模では効かなかった」が正確である。

\subsection{パラメータ数}

部分空間を変えて\emph{同一のレシピ}で当てはめた結果
（\texttt{capacity\_series.py}、$n = 637$、固定validation $152$件）：
\begin{center}
\begin{tabular}{lcc}
\toprule
モデル & 自由度 & Top-1 \\
\midrule
公開表（当てはめなし） & $0$ & $38.8\%$ \\
$\alpha$, clash, $\beta$のみ & $5$ & $40.1\%$ \\
\textbf{sym}（対称零和成分） & $\mathbf{12}$ & $\mathbf{50.0\%}$ \\
add（行＋列成分） & $23$ & $48.0\%$ \\
full & $144$ & $46.1\%$ \\
複合体を大きさで2分割 & $288$ & $47.4\%$ \\
同 4分割 & $576$ & $46.1\%$ \\
\bottomrule
\end{tabular}
\end{center}
\textbf{$12$個の自由度が$144$個を上回る}。パラメータを増やす方向は効かない。

これとは別に、実際に\emph{配備している}full-144の表（別レシピで当てはめたもの）は
同じvalidationで$42.76\%$であり、同条件で当てはめたsymの$50.00\%$は
これを$+7.24$\,pp 上回る。二つの比較は独立なので、
$42.76$と上の表を同じ列に並べてはいけない。

\textbf{bucket の失敗は「原子タイプを細分するな」を意味しない}（この区別は重要である）。
bucketは複合体を大きさで層別して\emph{同じ特徴に別の重み}を与えるものであり、
原子タイプの細分は\emph{いま同一視している接触を区別する}ものである。
しかもbucket2は$n = 637$でもまだ上を向いている（$34.9 \to 47.4$）ので、
飽和したとも言えない。前者の失敗は後者を否定しない。

ここで sym 部分空間とは、$e = e_0 + B\theta$ と書いたとき
$B$の列が「全体の水準」と「対称な零和モード」
$\frac{1}{\sqrt{24}}(a\otimes\mathbf{1} + \mathbf{1}\otimes a)$
（$a$は零和ベクトル）で張られるものである。

\subsection{その他}

\begin{itemize}
  \item hard negative採掘（round 1）：null。
        新規負例の$69.3\%$が重複、最難関より$1.87$\,SD下。
        \emph{射程}：seed $0$のみ、正例とPSCを固定、$1$ラウンド、固定プール。
        採掘あり対なしは$1$勝$0$敗（$p = 1$）で、end-to-endでは測っていない。
  \item 探索深度$K$の拡大：救えるのは最大$13/250 = 5.2\%$。
  \item 精密化後の分布で表を学習し直す：$123$件中$+1$件（$+0.82$\,pp）。
  \item basin安定性（終点を叩いて戻る割合）：単独ではAUC $0.738$の判別力を持つが、
        スコアと組み合わせると寄与ゼロ。学習を挟むと$-0.25$--$-4.28$\,pp。
        \emph{射程}：複合体内AUCであり、難しい訓練分布・上位$10$クラスタのみ・
        分類損失でのnullである。
\end{itemize}

これらはいずれも\textbf{測った条件の中でのnull}であって、
一般的な否定ではない。射程を明記したのはそのためである。

\section{結果と現在地}

公式PINDERベンチマーク（\texttt{pinder\_s}/holo、$250$件、$\mathrm{Max}(\mathrm{Top}\,1)$）：

\begin{center}
\begin{tabular}{llcccc}
\toprule
 & 精密化 & Acceptable & Medium & High & $\dockq$中央値\\
\midrule
公開表（出発点） & なし        & $66.40$ & $61.60$ & $30.80$ & $0.71$\\
公開表 + 精密化   & 上位$5$件   & $74.40$ & $72.40$ & $60.40$ & $0.88$\\
学習（凸）のみ    & なし        & $76.80$ & $72.40$ & $38.00$ & $0.76$\\
学習 + 精密化     & 上位$5$件   & $80.40$ & $78.40$ & $68.40$ & $0.91$\\
\textbf{学習 + 精密化} & \textbf{上位$\mathbf{50}$件}
                  & $\mathbf{84.00}$ & $\mathbf{82.40}$ & $\mathbf{73.60}$ & $\mathbf{0.91}$\\
\midrule
PatchDock & --- & $53.6$ & $52.4$ & $43.6$ & ---\\
FRODOCK   & --- & $89.6$ & $87.6$ & $84.8$ & ---\\
HDOCK     & --- & $94.4$ & $94.0$ & $93.6$ & ---\\
\bottomrule
\end{tabular}
\end{center}

\textbf{精密化する候補数$K$を揃えること}が比較の前提である。
$K = 5$と$K = 50$は同じ列に見えるが別の実験であり、
最終行だけが$K = 50$である。

$2\times2$の分解（公開表／学習表 $\times$ 精密化なし／あり）は
\textbf{$K = 5$の$4$条件のみ}に適用できる。それによれば、
Highの伸びは\textbf{ほぼ全部が精密化}であり
（精密化 $+29.6$/$+30.4$\,pp、学習 $+7.2$/$+8.0$\,pp、交互作用 $+0.8$）、
Acceptableでは学習の寄与が大きいが交互作用が$-4.4$\,ppと負である
（両者が同じ複合体を部分的に直している）。
すなわち\textbf{今回の伸びの大半は学習なしでも得られた}。
本研究の学習が寄与しているのは精密化の上に$+6.0$--$8.0$\,ppである。

\section{未解決の問題}

\begin{enumerate}
  \item \textbf{順位づけの$10.8$\,pp。}
        精密化後のプールに同じ凸学習を当てはめても$+0.82$\,ppに留まり、
        \emph{試した共有表のレシピでは}回収できなかった
        （共有表のTop-1の真の最大値は未確定である。$n \ge 25$でMILPが破綻したため）。
        非加法なモデル、あるいは現在のスコアが見ていない特徴量が候補になる。
        basin安定性は失敗した。
  \item \textbf{Highの品質。} Oracle $79.60\%$がHDOCKの$93.6\%$に届かない。
        「精密化エネルギー」（正解で最小になる滑らかな別関数）の学習が候補だが、
        その前に\emph{探索範囲が原因か、エネルギーが原因か}を切り分ける必要がある。
  \item \textbf{sym(12)の実地検証。} 固定プールで$+7.24$\,pp だが、
        表を変えると探索が返す候補も変わるため、end-to-end評価が必要。
  \item \textbf{方法論上の限界。} 本ベンチマークを設計判断に反復使用しているため、
        現在の数値は探索的である。未使用のcohort（\texttt{pinder\_xl} $\setminus$
        \texttt{pinder\_s}、$1705$件、クラスタ重複$0$を実測）での確認が要る。
  \item \textbf{holo redocking限定。} すべて結合型どうしの再ドッキングであり、
        PINDERの3設定で最も易しい。apo（非結合型）や予測構造を入力にしたときの
        外的妥当性は測っていない。
  \item \textbf{固定プールとend-to-endの対応。} 容量・部分空間の診断はすべて
        \emph{固定した}候補プール上の再順位づけで得たものである。
        表を変えると探索が返す候補自体が変わるので、
        固定プールでの順位と実際の探索結果の順位が一致する保証はない。
\end{enumerate}

\section{方法論としての要点}

本研究の技術的な核は、式\eqref{eq:affine}のアフィン性である。これにより
\begin{itemize}
  \item 学習が\textbf{凸}になるので、局所解に捕まる心配がなく、
        性能が出ないときに「最適化のせい」を実質的に排除できる；
  \item 複合体ごとの\textbf{分離可能性}を線形計画で、
        共有表の\textbf{達成値}を凸代理問題で測れるので、
        「容量のせい」かどうかを診断できる；
  \item 特徴量をキャッシュしておけば任意のパラメータで即座に再スコアできるので、
        診断が安価である。
\end{itemize}

一方、この性質は「$144$個のパラメータの線形モデル」という
古典的な設定に依存しており、ニューラルネットワークに基づく手法には移せない。
逆に、FFTに乗る相関分解は離散化した探索空間を\emph{網羅的に}評価できるので、
候補回収率は$96$--$97\%$と高い
（サンプリングに基づく手法にはこの保証がない）。

結論を正確に述べるとこうなる。
\textbf{FFT互換な相関分解は、離散探索空間を網羅的に評価できるという強みを持つ。
一方、本研究で用いた「$12$原子タイプ・単一の共有表・線形な接触特徴」という
具体的な構成は、再順位づけにおいて限界を示した。}
これはペア加法スコア\emph{一般}の限界を証明するものではない。
より細かい原子タイプ、複数の距離シェル、複合体依存の表といった
——いずれもFFTに乗る——拡張はまだ試されていない。

\end{document}
